
\subsection*{6. Appendice~: \'Etude des @@M0@@-faisceaux sur @@M1@@.}

\textbf{6.1.} Soient @@M2@@ un corps alg\'ebriquement clos, et @@M3@@. Soient @@M4@@ un faisceau coh\'erent sur @@M5@@, \`a polyn\^ome de Hilbert
@@M6@@
et
@@M7@@
une suite d'entiers @@M8@@. Soient
@@M9@@
Enfin, pour tout entier @@M10@@, d\'esignons le plus grand faisceau coh\'erent de @@M11@@ \`a support de dimension @@M12@@, par @@M13@@, et soit @@M14@@.

\textbf{Lemme 6.2.} \emph{Sous les conditions de 6.1, consid\'erons une suite exacte}
@@M15@@
\emph{avec @@M16@@. Alors @@M17@@ est @@M18@@-r\'egulier avec @@M19@@ o\`u @@M20@@ est le @@M21@@-i\`eme @@M22@@-polyn\^ome (1.9).}

En effet, @@M23@@ est \'evidemment un @@M24@@-faisceau (@@M25@@ termes), donc @@M26@@ est un @@M27@@-faisceau d'apr\`es 1.6. Gr\^ace \`a 2.10, @@M28@@. Donc l'assertion r\'esulte de 1.11 et 1.10.

\textbf{Proposition 6.3.} \emph{Sous les conditions de 6.1, on suppose @@M29@@ un @@M30@@-faisceau. Alors on peut calculer par le moyen de polyn\^omes universels en les @@M31@@, @@M32@@ et @@M33@@, des entiers @@M34@@ et @@M35@@ tels qu'on ait une suite exacte}
@@M36@@
\emph{De plus, @@M37@@ ne figure que dans le calcul de @@M38@@.}

